\section{模型与数值方法}

\subsection{五聚体 Heisenberg 链}

每个单胞包含五个自旋 $\frac{1}{2}$。以 $n$ 表示单胞编号，
$a=1,\ldots,5$ 表示胞内位置，哈密顿量写为
\begin{align}
H={}&J_1\sum_n
\left(
\bm S_{n,2}\cdot\bm S_{n,3}
+\bm S_{n,3}\cdot\bm S_{n,4}
\right)
\nonumber\\
&+J_2\sum_n
\left(
\bm S_{n,1}\cdot\bm S_{n,2}
+\bm S_{n,4}\cdot\bm S_{n,5}
\right)
\nonumber\\
&+J_3\sum_n
\bm S_{n,5}\cdot\bm S_{n+1,1}.
\label{eq:hamiltonian}
\end{align}
本文以 $J_1=1$ 为能量单位，并定义
\begin{equation}
\alpha=\frac{J_2}{J_1},
\qquad
\beta=\frac{J_3}{J_1}.
\label{eq:ratios}
\end{equation}

\subsection{孤立五聚体精确对角化}

孤立单胞的 Hilbert 空间维数为 $2^5=32$。我们对哈密顿量进行精确对角化，
并在每个简并能量子空间内进一步对角化总自旋算符
\begin{equation}
\bm S_{\mathrm{tot}}^2=
\left(\sum_{a=1}^{5}\bm S_a\right)^2,
\end{equation}
从而可靠地区分 $S=\frac{1}{2}$、$\frac{3}{2}$ 和 $\frac{5}{2}$
多重态。

为分析单自旋探针的局域激发通道，我们在固定能量、总自旋和五聚体镜像宇称的
完整简并子空间上构造投影器。对
$S^z(q)=5^{-1/2}\sum_{j=1}^{5}e^{-iqj}S_j^z$，谱权定义为
\begin{equation}
W_\lambda(q)=\frac{1}{d_g}
\operatorname{Tr}\left(P_g S^z(-q)P_\lambda S^z(q)\right),
\label{eq:local-projector-weight}
\end{equation}
其中 $P_g$ 是基态双重态投影器，$P_\lambda$ 是目标多重态投影器。
这一定义避免了简并本征矢在数值对角化中的任意旋转，并允许直接检查动量和规则、
$q=0$ 守恒以及镜像宇称选择定则。

\subsection{低能投影与有限链 DMRG}

在孤立五聚体基态双重态张成的子空间内，定义有效自旋
$\bm T_n$，则各站点自旋满足
\begin{equation}
P\bm S_{n,a}P=c_a\bm T_n.
\end{equation}
胞间耦合的一阶有效哈密顿量为
\begin{equation}
H_{\mathrm{eff}}=NE_5+J_{\mathrm{eff}}
\sum_{n=1}^{N-1}\bm T_n\cdot\bm T_{n+1}+O(\beta^2),
\qquad J_{\mathrm{eff}}=\beta c_1c_5.
\label{eq:effective-chain}
\end{equation}
投影系数由完整基态双重态计算，并通过两五聚体 ED 检查一阶能级劈裂与
$O(\beta^2)$ 剩余项。

有限链采用开放边界条件，总格点数为 $L=5N$。基态通过守恒总
$S^z$ 的矩阵乘积态 DMRG 得到。代表参数点使用 N\'{e}el、平移 N\'{e}el 和
双畴壁产品态独立初始化；只有能量密度、方差和目标自旋扇区同时通过门槛的最低
能态进入分析。相邻 $\beta$ 点可使用已收敛 MPS 加速，但不把 warm start 当作
独立收敛证据。

当前主要计算使用截断阈值 $10^{-10}$、最大键维数最高 $1500$，并对
$L=40,60,80$ 进行有限尺寸比较。不同观测量的详细参数应在正式投稿前整理为
补充材料表格。

\subsection{关联函数与弦关联}

普通连接自旋关联定义为
\begin{equation}
C(i,j)=
3\left[
\langle S_i^zS_j^z\rangle
-\langle S_i^z\rangle\langle S_j^z\rangle
\right].
\end{equation}
五聚体包络通过每五个格点提取一次局域极值，并与幂律
$|C(r)|\propto r^{-\eta}$ 和指数形式比较。

弦关联采用跨越五聚体单胞边界的端点约定。奇数位胞距离上的交错包络记为
\begin{equation}
\widetilde{\Ostr}(d)=
(-1)^{(d-1)/2}\Ostr(d),
\qquad d=1,3,5,\ldots
\end{equation}
弦关联拟合只使用 $3\le d\le 2N/3$ 的体区窗口，以排除短距离结构和
OBC 最长距离处的边缘效应。弦关联在本文中作为正区诊断量之一，不单独用于命名相。

\subsection{纠缠与动力学结构因子}

纠缠性质通过不同切割位置的 Schmidt 权重和纠缠熵分析。
动力学结构因子定义为
\begin{equation}
S^{zz}(q,\omega)=-\frac{1}{\pi}\operatorname{Im}
\left\langle b_q\left|
\frac{1}{\omega+E_0+\mathrm{i}\eta-H}
\right|b_q\right\rangle,
\qquad
|b_q\rangle=S_q^z|\psi_0\rangle,
\end{equation}
其中
$S_q^z=L^{-1/2}\sum_{j=1}^{L}\exp(\mathrm{i}qj)S_j^z$。
为对齐参考论文，主谱采用 Krylov 修正向量 DMRG 计算，展宽固定为
$\eta=0.1J_1$；实时 TEBD 仅作为独立方法交叉检查。

在进入生产尺寸前，我们在 $L=10$、总 $S^z=0$ 扇区将修正向量谱与完整
精确谱逐点比较。对 $q=\pi$ 和 $0\leq\omega/J_1\leq3$ 的 61 个频率点，
两种方法的谱函数相对 $L_2$ 误差为 $8.5\times10^{-12}$，积分权重相对误差为
$7.8\times10^{-13}$，最大线性方程相对残差为 $5.8\times10^{-7}$。

随后在 $L=20$、$q=\pi$ 和 $\omega=0.2J_1$ 进行长度与截断标定。
当修正向量截断阈值依次取 $10^{-10}$、$10^{-12}$ 和 $10^{-14}$ 时，
线性方程相对残差从 $3.80\times10^{-4}$ 单调下降到
$4.67\times10^{-5}$ 和 $5.62\times10^{-6}$，而最紧两档的谱强度相对差仅为
$2.33\times10^{-8}$。因此，较大尺寸计算将基态 DMRG 截断阈值与修正向量
截断阈值分别记录；当前修正向量基线取 $10^{-14}$。

相同参数下的 $L=40$ 单点探针得到相对残差 $7.48\times10^{-6}$、实际修正
向量键维数 256 和 $1.50$ GiB 峰值内存。修正向量求解时间相对 $L=20$ 增长
16.81 倍，对应本次两点长度估计的经验时间指数约为 4.07。一步外推给出的
$L=80$ 困难冷启动频率点约需 11.84 小时，因此在提交频率网格前必须先验证
相邻频率 warm start 和跨进程 checkpoint 恢复。

生产计算采用 $L=80$ 作为与参考论文 16 个单胞相同的五聚体对应尺寸，并对
最大键维数、频率网格和独立 TEBD 结果进行收敛检查。
